% Generated draft: requires mathematical and source review.
\documentclass[11pt]{article}
\usepackage[margin=25mm]{geometry}
\usepackage{fontspec}
\setmainfont{Latin Modern Roman}
\usepackage{amsmath,amssymb,mathtools,unicode-math}
\setmathfont{Latin Modern Math}
\usepackage{enumitem,needspace,xurl,hyperref}
\setlength{\parindent}{0pt}
\setlength{\parskip}{5pt}
\setlength{\emergencystretch}{3em}
\newcommand{\N}{\mathbb{N}}
\newcommand{\Z}{\mathbb{Z}}
\newcommand{\Q}{\mathbb{Q}}
\newcommand{\R}{\mathbb{R}}
\newcommand{\C}{\mathbb{C}}
\newcommand{\F}{\mathbb{F}}
\newcommand{\A}{\mathbb{A}}
\newcommand{\PP}{\mathbb{P}}
\newcommand{\E}{\mathbb{E}}
\newcommand{\eps}{\varepsilon}
\newcommand{\dd}{\,\mathrm{d}}
\begin{document}
\section*{UnsolvedMath ID: 10900002}
\textbf{Source problem code:} catalogue entry
\subsection{Definitions and mathematical statement}
Let \(M\) be a hyperbolic \(3\)-manifold admitting a CR structure, not assumed spherical. In real dimension \(3\), a hypersurface-type CR structure consists of a rank-two distribution \(H\subset TM\) and a complex structure \(J:H\to H\), with the relevant integrability condition; no Levi nondegeneracy or strict pseudoconvexity hypothesis is printed. Spherical means locally CR equivalent to the standard CR \(3\)-sphere, and is expressly not required.

A real projective structure is an atlas into \(\R\PP^3\) with transition maps in \(\operatorname{PGL}(4,\R)\). A convex projective structure is described by suitable developing and holonomy data with convex developing image; proper versus strict convexity and global developing-map requirements must be fixed. Let \(\mathcal D_{\mathrm{conv}}(M)\) denote the corresponding deformation space with the intended marking/equivalence conventions.

Choi asks whether the deformation theory of \(\mathcal D_{\mathrm{conv}}(M)\) can be understood in terms of the CR structure on \(M\). This asks for a precise construction or description relating the CR data to the projective deformations, not a supplied equality of two already defined moduli spaces. The source gives no map, parameterization, infinitesimal comparison or obstruction statement to prove. It proposes hyperbolic Coxeter \(3\)-orbifolds as a potentially easier setting for the same investigation. Completeness, compactness, volume and end/boundary conditions on \(M\), as well as what CR data vary or remain fixed, are not specified.
\subsection{Short English statement}
Can a CR structure, even a nonspherical one, explain convex real projective deformations of a hyperbolic \(3\)-manifold? Investigate hyperbolic Coxeter \(3\)-orbifolds as the suggested easier case.
\subsection{Sources}
\begin{itemize}
\item Choi, attributed item 1.2 in Delp, Hoffoss and Manning, \emph{Problems In Groups, Geometry, and Three-Manifolds} (2015), supplied PDF page 1. \url{https://arxiv.org/pdf/1512.04620#page=1}; source-status link \url{https://arxiv.org/abs/1512.04620}. Recorded access 2026-07-29; extraction source-tex checked against rendered PDF text, as reported by the supplied record.
\item Local UnsolvedMath registry, \texttt{data/unsolvedmath/problems.json}, permanent ID \(10900002\), exact source code \texttt{catalogue entry}. Its OPEN-TRIAGE summary mentions established surrounding holonomy openness/closedness work but supplies no exact additional reference or verified resolution of this question.
\end{itemize}
\end{document}
